% Source: https://artofproblemsolving.com/community/c3h1760794_square_root_inequality_verification
% Stable identity: aops-topic-1760794-q01
The square root requires $5-2x\ge0$, i.e.\ $x\le\dfrac{5}{2}$.

\textbf{Case $x<1$.} The right side is negative and the left side is nonnegative, so every $x<1$ is a solution (all such $x$ satisfy $x\le\dfrac{5}{2}$).

\textbf{Case $1\le x\le\dfrac{5}{2}$.} Both sides are nonnegative, so the inequality is equivalent to $5-2x>(x-1)^2=x^2-2x+1$, i.e.\ $x^2<4$, i.e.\ $-2<x<2$. Together with $1\le x\le\dfrac{5}{2}$ this gives $1\le x<2$.

The solution set is the \emph{union} of the two cases: $x<1$ or $1\le x<2$, that is, $\boxed{x<2}$. The answer key is right.

The student's case work is correct, but at the end the condition $-2<x<2$ from the second case was intersected with everything (including the first case's $x<1$) instead of being combined only with $1\le x\le\dfrac{5}{2}$ and then joined to the first case. For example $x=-5$ is a solution, since $\sqrt{15}>-6$, but it is excluded by $-2<x<2$.
